% TOP_PROBLEM: {"id": 30001063, "title": "Falconer Distance Conjecture", "category_id": 6, "category": {"id": 6, "name": "geometry", "display_name": "Geometry", "description": "Euclidean and non-Euclidean geometry, geometric structures.", "slug": "geometry", "order_index": 6, "created_at": "2026-07-31T15:26:25.671Z"}, "status": "open"}
% FIRST_POSTED: 2026-09-27
% Standalone research entry 209; batch 208--217.
% Original section material preserved from Pasted text(3).txt.
% Research additions: 27 September 2026.
% Compile twice with: lualatex 30001063.tex
% !TeX program = lualatex
% Compile twice: lualatex open_problems_500.tex
% All references and prose are embedded; no BibTeX database or external inputs.
\documentclass[11pt,a4paper]{article}
\usepackage[margin=24mm,headheight=14pt]{geometry}
\usepackage{fontspec}
\setmainfont{Latin Modern Roman}
\setsansfont{Latin Modern Sans}
\setmonofont{Latin Modern Mono}
\usepackage{amsmath,amssymb,mathtools}
\usepackage{unicode-math}
\setmathfont{Latin Modern Math}
\usepackage{microtype}
\usepackage{enumitem,xurl,needspace}
\usepackage[dvipsnames]{xcolor}
\usepackage{hyperref}
\hypersetup{unicode=true,colorlinks=true,linkcolor=MidnightBlue,urlcolor=MidnightBlue,
 pdftitle={500 Mathematical Problem Statements},pdfauthor={Source-annotated compilation},bookmarksnumbered=true}
\usepackage{bookmark}
\usepackage{tocloft}
\setlength{\cftsecnumwidth}{3em}
\cftsetpnumwidth{2.5em}
\cftsetrmarg{4em}
\usepackage{titlesec}
\titleformat{\section}{\Large\bfseries\raggedright}{\thesection}{0.7em}{}
\usepackage{fancyhdr}
\pagestyle{fancy}\fancyhf{}
\fancyhead[L]{\small\sffamily Mathematical problem statements}
\fancyhead[R]{\small Source edition 22}
\fancyfoot[C]{\thepage}
\setlength{\parindent}{0pt}
\setlength{\parskip}{5pt plus 1pt}
\setlength{\emergencystretch}{3em}
\setcounter{tocdepth}{1}
\setcounter{secnumdepth}{1}
\setlist[itemize]{leftmargin=3em,itemsep=5pt,topsep=4pt,parsep=0pt}
\allowdisplaybreaks
\newcommand{\N}{\mathbb{N}}
\newcommand{\Z}{\mathbb{Z}}
\newcommand{\Q}{\mathbb{Q}}
\newcommand{\R}{\mathbb{R}}
\newcommand{\C}{\mathbb{C}}
\newcommand{\F}{\mathbb{F}}
\newcommand{\A}{\mathbb{A}}
\newcommand{\PP}{\mathbb P}
\newcommand{\E}{\mathbb{E}}
\newcommand{\eps}{\varepsilon}
\newcommand{\dd}{\,\mathrm d}
\newcommand{\sref}[2]{\hyperlink{source:#1:#2}{[S#2]}}
\newcommand{\eref}[2]{\hyperlink{extra:#1:#2}{[E#2]}}
% Turn the next switch off for a shorter reading copy. The quotations preserve
% every exactTarget field, including qualifications omitted from a short summary.
\newif\ifshowcatalogue
\showcataloguetrue
\newcommand{\cataloguescope}[1]{\ifshowcatalogue
 \paragraph{Exact scope in the supplied catalogue.}
 \begin{quote}\small #1\end{quote}\fi}

% Standalone wrapper; no external inputs or bibliography files are required.
\fancyhead[L]{\small\sffamily Research notebook: section 30001063}
\fancyhead[R]{\small 27 September 2026}
\hypersetup{pdftitle={Research attempt 209},pdfauthor={Research notebook}}
\setlength{\parskip}{4pt plus 1pt}
\setlist[itemize]{before=\small\interlinepenalty10000,itemsep=3pt}
\begin{document}
\setcounter{section}{0}
% ================================================================
% problemId: 30001063
% Canonical problem ID: 30001063
% exactTarget SHA-256: 921d62550cbdd9d0a27caf2ef317882b8ac32f14da90ce3a02ddd79f9ec3760e
\clearpage
\section*{\texorpdfstring{Falconer Distance Conjecture}{Falconer Distance Conjecture}}\label{30001063}
\subsection{Definitions and mathematical statement}
For compact $E\subset\R^d$, define its Euclidean distance set
$\Delta(E)=\{|x-y|:x,y\in E\}\subset[0,\infty)$. Hausdorff dimension is
$\dim_H E=\inf\{s:\mathcal H^s(E)=0\}$, with $\mathcal H^s$ defined by arbitrarily fine covers and sums of $s$th powers of diameters. The conjecture is
\[
 \dim_H E>d/2\quad\Longrightarrow\quad \mathcal L^1(\Delta(E))>0
 \qquad(d\ge2).
\]
The threshold inequality is strict. Positive one-dimensional measure is stronger than merely uncountably many distances or a distance set of Hausdorff dimension one. The set need not carry positive $d$-dimensional volume.
\par\smallskip\textit{Formulation and concept references:} \sref{30001063}{1}.
\cataloguescope{Prove or disprove that for every integer d>=2 and every compact set E in R\textasciicircum{}d with Hausdorff dimension greater than d/2, the distance set Delta(E)=\{|x-y|: x,y in E\} has positive one-dimensional Lebesgue measure.}
\subsection{Short English statement}
Must a compact Euclidean set whose Hausdorff dimension exceeds half the ambient dimension determine a positive-measure set of distances?
% BEGIN RESEARCH ADDITION 209
\subsection{Research attempt}

\paragraph{Current frontier.}
The planar theorem of Guth--Iosevich--Ou--Wang reaches dimension $>5/4$,
with a pinned positive-measure conclusion.\eref{30001063}{1}
Du--Ou--Ren--Zhang, Theorem 1.1 in the 2024 revision, obtain the threshold
$d/2+1/4-1/(8d+4)$ for $d\ge3$, again with a pinned conclusion.
\eref{30001063}{2} These exceed $d/2$. Their use of carefully selected or modified
measures is important: a conjectural estimate for \emph{every} measure on
$E$ would be stronger than the set-existence statement. The attempt below
only asks to find one suitable probability measure.

\paragraph{Attack route.}
Fourier $L^2$ estimates give a strong way to force positive measure, but may
pay too much for rare, highly populated distance annuli. I investigate an
entropy budget instead. It penalizes annular imbalance logarithmically and
has an exact nonnegative increment at each dyadic refinement. The possible
advantage is that one can tolerate some large collision energies; the
missing lemma is summability of the actual distance-measure imbalance,
not merely a dimension-one conclusion. A sparse-digit construction tests
whether a weaker asymptotic entropy statement could suffice.

\paragraph{Attempt.}
\textbf{1. Normalize the distance measure without changing the target.}
Fix compact $E$ with $\dim_H E>d/2$ and choose $L>\operatorname{diam}(E)$.
For a Borel probability measure $\mu$ supported on $E$, let
\[
 \nu(A)=(\mu\times\mu)\{(x,y): |x-y|/L\in A\},\qquad A\subset[0,1).
\]
Its support lies in $\Delta(E)/L$, a compact set. The measure need not have
full support on $E$; any positive-measure subset of the distance set suffices.
Let $I_{j,k}=[k2^{-j},(k+1)2^{-j})$, $0\le k<2^j$, and define
\[
 p_{j,k}=\nu(I_{j,k}),\qquad
 H_j=-\sum_k p_{j,k}\log p_{j,k},\qquad
 D_j=j\log2-H_j,
\]
with $0\log0=0$ and natural logarithms. The normalization $L>\operatorname{diam}(E)$
avoids an endpoint mass at $1$.

\textbf{2. Bounded entropy deficit forces positive measure.}
For completeness, a direct support argument avoids assuming a density for
$\nu$. If $N_j$ cells have positive mass, concavity of the logarithm gives
$H_j\le\log N_j$, and hence
\begin{equation}\label{30001063:support}
 2^{-j}N_j\ge e^{-D_j}.
\end{equation}
Let $U_j$ be the union of the \emph{closed} cells corresponding to positive
masses. The sets $U_j$ are compact and nested: a positive-mass child has a
positive-mass parent. Their intersection is contained in $\operatorname{supp}\nu$.
Indeed, a point outside that support has a neighborhood of zero mass and
cannot belong to such a cell once the cell diameter is sufficiently small.
Closing cells changes no Lebesgue measure. If $D_j\le C$ for every $j$,
continuity of Lebesgue measure from above and \eqref{30001063:support} yield
\begin{equation}\label{30001063:positive}
 |\Delta(E)|\ge L|\operatorname{supp}\nu|\ge Le^{-C}>0.
\end{equation}
There is no assumption here that Hausdorff dimension one already implies
positive measure.

\textbf{3. An exact annular-imbalance identity.}
For a parent of positive mass, put
\[
 \beta_{j,k}=\frac{p_{j+1,2k}-p_{j+1,2k+1}}{p_{j,k}}\in[-1,1],
 \qquad
 \Phi(b)=\frac{(1+b)\log(1+b)+(1-b)\log(1-b)}2.
\]
Set the summand below to zero when $p_{j,k}=0$. Expanding the entropy of
both children gives the exact chain rule
\begin{equation}\label{30001063:chain}
 D_{j+1}-D_j=\sum_k p_{j,k}\Phi(\beta_{j,k})=:B_j\ge0,
 \qquad D_0=0.
\end{equation}
Here $\Phi(0)=0$, $\Phi(\pm1)=\log2$, and
$\Phi''(b)=1/(1-b^2)$ on $(-1,1)$.
Thus $\Phi(b)\ge b^2/2$, while $\Phi(b)\le(2/3)b^2$ for $|b|\le1/2$.
All endpoint values are limits, so a parent whose entire mass chooses one
child is included rather than silently excluded.

This proves the following sufficient criterion, including its constants.
Define
\[
 a_j=\sum_{k:\,|\beta_{j,k}|>1/2}p_{j,k},\qquad
 b_j=\sum_{k:\,|\beta_{j,k}|\le1/2}p_{j,k}\beta_{j,k}^2.
\]
If $\sum_j a_j<\infty$ and $\sum_j b_j<\infty$, then
\begin{equation}\label{30001063:budget}
 |\Delta(E)|\ge
 L\exp\left(-\log2\sum_j a_j-\frac23\sum_j b_j\right)>0.
\end{equation}
The proof is just the upper bounds for $\Phi$, telescoping
\eqref{30001063:chain}, and \eqref{30001063:positive}. This allows very biased cells
provided their total probability mass, summed over scales, is finite.

\textbf{4. What the missing geometric estimate actually says.}
In terms of the original points, define the alternating annular indicator
\[
 A_{j,k}(t)=\mathbf1_{I_{j+1,2k}}(t)-\mathbf1_{I_{j+1,2k+1}}(t).
\]
Then the difference tested in the imbalance is exactly
\[
 p_{j+1,2k}-p_{j+1,2k+1}
   =\iint_E A_{j,k}(|x-y|/L)\,d\mu(x)d\mu(y).
\]
A completion would establish, for every admissible $E$, the existence of
$\mu$ for which the sum of the $B_j$ in \eqref{30001063:chain} is finite.
This is a sufficient condition, not an asserted equivalent reformulation
of Falconer: a positive-measure distance set need not be witnessed by a
measure with this entropy bound.

For comparison let $\mathcal E_j=2^j\sum_kp_{j,k}^2$. Jensen's inequality
with weights $p_{j,k}$ gives
\[
 D_j=\sum_kp_{j,k}\log(2^jp_{j,k})\le\log\mathcal E_j.
\]
Uniform dyadic $L^2$ control therefore implies the entropy criterion. The
converse is false even for ordinary densities: an integrable probability
density proportional near zero to
$x^{-1/2}(\log(e/x))^{-1/4}$ has finite $\int f\log f$ but infinite
$\int f^2$. Its dyadic deficits are bounded by $\int f\log f$ through
cellwise Jensen. This demonstrates a real analytic weakening, although
it does not establish that a geometric construction can exploit it.

\textbf{5. A sharp warning about sublinear losses.}
Take independent binary digits $\xi_r$, except that $\xi_r=0$ whenever
$r$ is a perfect square, and put
\[
 X=\sum_{r\ge1}\xi_r2^{-r};
\]
all other digits are fair. Let $F$ be the compact set of all such expansions
and $\sigma$ the resulting probability law. Ambiguous dyadic expansions
form a countable set of $\sigma$-measure zero, because there are infinitely
many free digits. At level $j$ exactly $2^{j-\lfloor\sqrt j\rfloor}$ cells
have positive mass, all equal. Consequently
\begin{equation}\label{30001063:sparse}
 H_j(\sigma)=(j-\lfloor\sqrt j\rfloor)\log2,
 \quad D_j(\sigma)=\lfloor\sqrt j\rfloor\log2,
 \quad \mathcal E_j(\sigma)=2^{\lfloor\sqrt j\rfloor}.
\end{equation}
These occupied cells cover $F$ up to their endpoints and have total length
tending to zero, so $|F|=0$. Nevertheless $\dim_H F=1$. To verify the lower
bound, for every $s<1$ the largest level-$j$ cell mass is at most
$C_s2^{-sj}$. An interval of length between $2^{-j-1}$ and $2^{-j}$ meets
at most two level-$j$ cells apart from null endpoints. Hence
$\sigma(I)\le C'_s|I|^s$, which implies $\dim_H F\ge s$ by applying this
inequality to arbitrary small covers. Let $s\uparrow1$.

This example has $D_j=o(j)$ and $\mathcal E_j=2^{o(j)}$, yet zero measure.
Its entropy cost is concentrated at the sparse forced-digit scales:
$B_{r-1}=\log2$ for square $r$ and zero otherwise. An arbitrarily long
run of balanced scales gives no control of the total future deficit.

\paragraph{Stress test.}
The sparse-digit example is a measure on the line, not a constructed
counterexample distance set of a high-dimensional $E$. Its role is to
falsify the inference from sublinear entropy loss to positive measure.
The compact-support proof handles dyadic boundaries explicitly and does
not interchange an unproved density limit with integration. The chain
rule is finite at each scale and telescopes monotonically, so its use
requires no cancellation of a conditionally convergent series.

The central unproved claim is the existence of a measure satisfying the
summable annular-imbalance budget for \emph{every} compact set of dimension
$>d/2$. Frostman ball bounds alone do not furnish this: annuli constrain
pairs and their alternating differences, not just mass inside balls.
Known pinned-measure modifications should not be assumed to preserve this
new budget without a separate proof.\eref{30001063}{1}\eref{30001063}{2}

\Needspace{11\baselineskip}
\paragraph{Outcome.}
\textbf{Outcome: Conditional reduction.}

A summable entropy-imbalance estimate implies the exact positive-measure
conclusion, with the explicit lower bound \eqref{30001063:budget}. The proof
also separates this estimate from both uniform $L^2$ control and a mere
full-dimensional distance conclusion. No summable budget has been proved
under the conjecture's hypotheses alone, so no new dimension threshold or
solution of Falconer is asserted. The criterion and stress tests are
not claimed to be new in the literature without further comparison.
% END RESEARCH ADDITION 209
\subsection{Sources}
\begin{itemize}\raggedright
\item[\textnormal{[S1]}]\hypertarget{source:30001063:1}{} Krystal Taylor and Samantha Sandberg-Clark, Triangles in the plane and arithmetic progressions in thick compact subsets of R\textasciicircum{}d, Introduction.
\url{https://www.cambridge.org/core/journals/canadian-journal-of-mathematics/article/triangles-in-the-plane-and-arithmetic-progressions-in-thick-compact-subsets-of-mathbb-rd/B5F9FC2F5A15CA0577BEDC96031B84CB}.
{\small\textit{Catalogue formulation source.}}
% BEGIN RESEARCH REFERENCES 209
\item[\textnormal{[E1]}]\hypertarget{extra:30001063:1}{} Larry Guth, Alex Iosevich, Yumeng Ou and Hong Wang,
\emph{On Falconer's distance set problem in the plane}, Inventiones Mathematicae
219 (2020), 779--830; arXiv:1808.09346, main theorem and the good/bad measure decomposition.
\url{https://arxiv.org/abs/1808.09346}.
{\small\textit{Consulted 27 September 2026.}}
\item[\textnormal{[E2]}]\hypertarget{extra:30001063:2}{} Xiumin Du, Yumeng Ou, Kevin Ren and Ruixiang Zhang,
\emph{New improvement to Falconer distance set problem in higher dimensions},
arXiv:2309.04103v2 (22 October 2024), Theorem 1.1 and Section 1.
\url{https://arxiv.org/abs/2309.04103}.
{\small\textit{Consulted 27 September 2026; the version used includes dimension three.}}
% END RESEARCH REFERENCES 209
\end{itemize}

\end{document}
